\chapter{Kontextlänge und Pricing -- die harten Fakten}

\label{chap:kontext_pricing}

% ============================================================
% LERNZIELE
% ============================================================
\begin{lernziele}
Nach diesem Kapitel kannst du:
\begin{itemize}
    \item Die Kostenstruktur von LLM-APIs verstehen und berechnen
    \item Kontextfenster effizient nutzen und Grenzen managen
    \item Die Trade-offs zwischen Kontextlänge, Kosten und Latenz bewerten
    \item Prompt-Caching-Strategien für wiederholte System-Prompts umsetzen
    \item Die Pricing-Modelle der wichtigsten Anbieter vergleichen
\end{itemize}
\end{lernziele}

% ============================================================
% MOTIVATION
% ============================================================
\section{Motivation}

Die meisten AI-Engineering-Fehler sind keine Architektur-Fehler -- sie sind
Kosten- und Kontext-Fehler. Ein Entwickler packt 50 Dokumente in einen
Prompt, weil das Modell 1M Token Kontext hat -- und wundert sich über eine
Rechnung von 5 € pro Query. Ein anderes Team nutzt GPT-4o für jede Anfrage,
obwohl GPT-4o-mini in 80\,\% der Fälle genauso gut ist.

Kontextlänge und Pricing sind die beiden harten Fakten, die jede
Architekturentscheidung bestimmen. Kontext kostet Geld -- jeder Token,
den du in den Prompt packst, wird berechnet. Und Kontext hat Grenzen:
Selbst bei 1M Token Kontextfenster findet das Modell nicht automatisch
die relevante Information.

In diesem Kapitel lernst du, Kosten zu kalkulieren, Kontext zu managen
und die richtige Balance zwischen Kontextlänge, Antwortqualität und
Preis zu finden.

% ============================================================
% GRUNDLAGEN
% ============================================================
\section{Grundlagen -- Kontext und Kosten verstehen}

\subsection{Das Context Window}

Das Context Window (Kontextfenster) ist die maximale Menge an Tokens,
die ein Modell in einem Durchlauf verarbeiten kann. Es umfasst:

\begin{itemize}[leftmargin=*]
    \item \textbf{System-Prompt}: Verhaltensdefinition, Regeln, Persona
    \item \textbf{User-Prompt}: Die aktuelle Anfrage
    \item \textbf{Chat-History}: Vorherige Nachrichten im Gespräch
    \item \textbf{RAG-Kontext}: Dokumente, die zur Beantwortung relevant sind
    \item \textbf{Output}: Die Antwort des Modells (auch Teil des Fensters)
\end{itemize}

Die Summe aller Komponenten darf das Context Window nicht überschreiten.

\begin{table}[H]
\centering
\begin{tabularx}{\linewidth}{lX}
\toprule
\textbf{Modell} & \textbf{Kontextfenster} \\
\midrule
GPT-4o / GPT-4o-mini & 128K Tokens \\
Claude 4 Sonnet / Opus & 200K Tokens \\
Gemini 2.5 Flash / Pro & 1M Tokens \\
Llama 3 (alle Größen) & 128K Tokens \\
Mistral Large & 128K Tokens \\
Qwen 2.5 & 128K Tokens \\
DeepSeek-V3 & 128K Tokens \\
\bottomrule
\end{tabularx}
\caption{Kontextfenster der wichtigsten Modelle}
\end{table}

\subsection{Input- vs. Output-Tokens}

APIs berechnen Input- und Output-Tokens getrennt:

\begin{description}[leftmargin=*,style=nextline]
    \item[Input-Tokens] Was du an das Modell schickst (Prompt + History)
    -- inklusive aller Context-Tokens. Typisch 70--90\,\% der Gesamtkosten
    bei RAG-Systemen.

    \item[Output-Tokens] Was das Modell generiert. Output-Tokens sind
    3--10x teurer als Input-Tokens, da sie rechenintensiver sind
    (autoregressive Generation).
\end{description}

\begin{verbatim}
    Gesamtkosten = Input-Tokens * Input-Preis
                 + Output-Tokens * Output-Preis
\end{verbatim}

\subsection{Token-Kosten pro Millionen Tokens}

\begin{table}[H]
\centering
\begin{tabularx}{\linewidth}{lrr}
\toprule
\textbf{Modell} & \textbf{Input (pro Mio. Tokens)} & \textbf{Output} \\
\midrule
GPT-4o-mini & 0,15 \$ & 0,60 \$ \\
GPT-4o & 2,50 \$ & 10,00 \$ \\
Claude 4 Sonnet & 3,00 \$ & 15,00 \$ \\
Claude 4 Opus & 15,00 \$ & 75,00 \$ \\
Gemini 2.5 Flash & 0,08 \$ & 0,30 \$ \\
Gemini 2.5 Pro & 1,25 \$ & 10,00 \$ \\
DeepSeek-V3 & 0,27 \$ & 1,10 \$ \\
Llama 3 (lokal) & Stromkosten & Stromkosten \\
\bottomrule
\end{tabularx}
\caption{API-Kosten pro Million Tokens}
\label{tab:api-kosten-kontext}
\end{table}

\begin{praxishinweis}
Der größte Kostenfehler: Jeden Request mit dem teuersten Modell zu bearbeiten. GPT-4o-mini reicht für 80\,\% aller Aufgaben. Ein Semantic Router, der einfache Queries automatisch an günstigere Modelle leitet, spart 60--80\,\% der Kosten.
\end{praxishinweis}

% ============================================================
% ARCHITEKTUR
% ============================================================
\section{Architektur -- Context-Management in Produktion}

Jeder LLM-Call durchläuft ein Context-Budgeting:

\begin{verbatim}
    Request kommt an
         |
         v
    [Context-Budget: max 100K Tokens]
         |
    +----+----+
    |         |
    RAG      Chat-History
    (50K)    (30K)
    |         |
    +----+----+
         |
    System-Prompt (5K)
         |
    User-Input (1K)
         |
    Summe: 86K < 100K => OK
\end{verbatim}

\subsection{Context-Window-Management-Strategien}

\begin{description}[leftmargin=*,style=nextline]
    \item[Sliding Window] Behalte nur die letzten N Nachrichten.
    Älteste Nachrichten werden verworfen. Einfach, effektiv, aber
    verliert Kontext bei langen Gesprächen.

    \item[Summary Compression] Fasse alte Nachrichten in einer
    Kurzzusammenfassung zusammen. Teurer (extra LLM-Call), aber
    erhält den Kontext.

    \item[Semantic Filtering] Behalte nur Nachrichten, die für die
    aktuelle Query relevant sind. Nutzt Embedding-Ähnlichkeit zur
    Selektion. Aufwändig, aber am besten für RAG-Systeme.

    \item[Hybrid] Sliding Window für kurze Konversationen + Summary
    Compression für lange Sessions. Der Standard.
\end{description}

% ============================================================
% THEORIE
% ============================================================
\section{Theorie -- Warum Kontext Grenzen hat}

\subsection{Lost in the Middle}

Das wichtigste Kontext-Phänomen: Modelle schenken dem Anfang und Ende
eines langen Prompts mehr Aufmerksamkeit als der Mitte. Studien zeigen,
dass die Retrieval-Genauigkeit in der Mitte eines 100K-Kontexts um
30--50\,\% niedriger sein kann als am Anfang oder Ende.

\begin{verbatim}
    Retrieval-Genauigkeit uber die Kontext-Position:
    
    100% |\             /|
         | \           / |
    80%  |  \         /  |
         |   \       /   |
    60%  |    \-----/    |
         | Anfang  Mitte  Ende
    40%  +-------------------+
             Kontext-Position
\end{verbatim}

Konsequenz: Auch bei grossen Kontextfenstern solltest du die wichtigsten
Informationen an den Anfang oder das Ende des Prompts setzen. Relevante
Dokumente gehören in die ersten 20\,\% oder letzten 10\,\% des Kontexts.

\subsection{Attention Scaling}

Mit wachsender Kontextlänge steigt die Rechenzeit quadratisch
(O(n²) durch die Self-Attention im Transformer). Ein Prompt mit
200K Tokens ist nicht 2x so teuer wie einer mit 100K Tokens --
sondern bis zu 4x teurer.

\begin{table}[H]
\centering
\begin{tabularx}{\linewidth}{lXX}
\toprule
\textbf{Kontextlänge} & \textbf{Rechenzeit} & \textbf{Relative Kosten} \\
\midrule
1K Tokens & 1x & 1x \\
10K Tokens & 10x & 10x \\
100K Tokens & 100x & 100x \\
200K Tokens & 400x & 400x (Attention: n²) \\
\bottomrule
\end{tabularx}
\caption{Skalierungsverhalten der Kontextverarbeitung}
\end{table}

\subsection{LLMs sind keine Suchmaschinen}

Ein häufiges Missverständnis: \glqq Das Modell hat 1M Token Kontext,
also werfe ich einfach alle meine Dokumente rein.\grqq{} Das funktioniert
nicht -- LLMs sind nicht dafür optimiert, in riesigen Textmengen die
Nadel im Heuhaufen zu finden. Dafür brauchst du RAG mit einem
Retrieval-System, das die relevanten Dokumente vorselektiert.

\begin{praxishinweis}
Lost in the Middle ist kein Randphänomen -- es betrifft jeden Prompt über 10K Tokens. Strukturiere deinen Kontext immer nach dem Prinzip: wichtigste Info an den Anfang, zweitwichtigste ans Ende, der Rest in die Mitte. Oder besser: nutze RAG, um die Dokumentenmasse vorab zu filtern.
\end{praxishinweis}

% ============================================================
% PRAXIS -- Kostenrechnung
% ============================================================
\section{Praxis -- Kostenkalkulation im Alltag}

\subsection{Beispielrechnung: RAG-System}

Ein RAG-Chatbot mit 10.000 täglichen Queries:

\begin{table}[H]
\centering
\begin{tabularx}{\linewidth}{lXr}
\toprule
\textbf{Position} & \textbf{Berechnung} & \textbf{Kosten/Tag} \\
\midrule
Input (System + RAG-Kontext) & 10K Tokens $\times$ 10.000 Queries & 250 \$ (GPT-4o) \\
Output (Antwort) & 500 Tokens $\times$ 10.000 Queries & 500 \$ (GPT-4o) \\
Gesamt GPT-4o & & 750 \$/Tag \\
Gesamt GPT-4o-mini & 10K $\times$ 10.000 $\times$ 0,15 \$/M & 15 \$/Tag \\
Gesamt Gemini Flash & 10K $\times$ 10.000 $\times$ 0,08 \$/M & 8 \$/Tag \\
\bottomrule
\end{tabularx}
\caption{Tägliche Kosten eines RAG-Chatbots mit 10.000 Queries}
\end{table}

Die Botschaft: GPT-4o-mini ist 50x günstiger als GPT-4o. Bei 10.000
täglichen Queries entscheidet das über 250.000 \$ Jahreskosten.

\subsection{Cache-Effekte nutzen}

Wenn dein System-Prompt konstant bleibt (z.B. 5K Tokens), kannst du
durch Prompt-Caching 50--90\,\% der Input-Kosten sparen:

\begin{verbatim}
    Ohne Caching:   10.000 Queries * 10K Input = 100 Mio Tokens/Tag
    Mit Caching:    10.000 Queries * 5K (dynamisch) + 5K (1x gecached)
                    = 50 Mio Tokens/Tag + 5K einmalig
    Ersparnis:      ~50% bei Input-Kosten
\end{verbatim}

% ===========================================================%
% CODE -- Token-Counting und Context-Budgeting
% ============================================================
\section{Codebeispiele}

\subsection{Token-Counting vor dem API-Call}

\begin{lstlisting}[language=Python, caption={Token-Counting mit tiktoken}, label=lst:token-counting]
import tiktoken

def count_tokens(text: str,
                 model: str = "gpt-4o") -> int:
    """Zahle die Tokens eines Textes fur ein bestimmtes Modell."""
    encoding = tiktoken.encoding_for_model(model)
    return len(encoding.encode(text))

def enforce_context_budget(
    messages: list[dict],
    max_tokens: int = 100_000,
    model: str = "gpt-4o",
) -> list[dict]:
    """Stelle sicher, dass der Prompt ins Context-Window passt."""
    total = 0
    budget_messages = []

    for msg in reversed(messages):
        tokens = count_tokens(msg["content"], model)
        if total + tokens > max_tokens:
            break
        budget_messages.insert(0, msg)
        total += tokens

    return budget_messages

# Beispiel: Chat-History managen
history = [
    {"role": "user", "content": "Frage 1"},
    {"role": "assistant", "content": "Antwort 1"},
    # ... viele weitere Nachrichten
    {"role": "user", "content": "Frage 50"},
]

budget = enforce_context_budget(
    history, max_tokens=50_000
)
print(f"Gekurzt von {len(history)} auf {len(budget)} Nachrichten")
\end{lstlisting}

\subsection{Kostenrechner für den Produktionseinsatz}

\begin{lstlisting}[language=Python, caption={Kostenrechner für LLM-API-Calls}, label=lst:cost-calculator]
from dataclasses import dataclass

@dataclass
class ModelPricing:
    input_price_per_million: float
    output_price_per_million: float

MODEL_PRICING = {
    "gpt-4o": ModelPricing(2.50, 10.00),
    "gpt-4o-mini": ModelPricing(0.15, 0.60),
    "claude-sonnet-4": ModelPricing(3.00, 15.00),
    "gemini-flash-2.5": ModelPricing(0.08, 0.30),
}

def estimate_cost(
    model: str,
    input_tokens: int,
    output_tokens: int,
    daily_queries: int = 1,
) -> dict:
    """Berechne Kosten fur einen oder mehrere API-Calls."""
    pricing = MODEL_PRICING[model]

    input_cost = (input_tokens / 1_000_000
                  * pricing.input_price_per_million)
    output_cost = (output_tokens / 1_000_000
                   * pricing.output_price_per_million)
    per_query = input_cost + output_cost

    return {
        "model": model,
        "cost_per_query": round(per_query, 5),
        "cost_per_1k_queries": round(per_query * 1000, 2),
        "cost_per_10k_queries": round(per_query * 10000, 2),
        "cost_monthly_30d": round(per_query * daily_queries * 30, 2),
    }

# Beispiel: Kosten fur ein RAG-System
costs = estimate_cost(
    "gpt-4o-mini",
    input_tokens=8_000,   # System + RAG-Kontext
    output_tokens=500,     # Antwort
    daily_queries=10_000,
)
print(f"Pro Query:   {costs['cost_per_query']} $")
print(f"Pro Monat:   {costs['cost_monthly_30d']} $")
\end{lstlisting}

\subsection{Prompt-Caching mit Anthropic}

\begin{lstlisting}[language=Python, caption={Prompt-Caching mit Cache-Control}, label=lst:prompt-caching]
# Prompt-Caching reduziert Input-Kosten um bis zu 90%
# fur statische Prompt-Teile (System-Prompt, Few-Shot)

response = client.chat.completions.create(
    model="claude-sonnet-4",
    messages=[
        {
            "role": "system",
            "content": [
                {
                    "type": "text",
                    "text": SYSTEM_PROMPT,
                    "cache_control": {
                        "type": "ephemeral"
                    },
                }
            ],
        },
        {
            "role": "user",
            "content": user_query,
        },
    ],
)

# Check Cache-Treffer
print(f"Cache Input: {response.usage.cache_read_input_tokens}")
print(f"Cache Create: {response.usage.cache_creation_input_tokens}")
\end{lstlisting}

\begin{praxishinweis}
Prompt-Caching funktioniert nur bei exaktem Prefix-Match. Stelle sicher, dass der System-Prompt immer am Anfang des Message-Arrays steht und sich zwischen Requests nicht ändert. Dynamische Elemente wie User-Name oder Datum gehören in den User-Prompt, nicht in den System-Prompt.
\end{praxishinweis}

% ===========================================================%
% BEST PRACTICES
% ============================================================
\section{Best Practices}

\begin{itemize}[leftmargin=*]
    \item \textbf{Starte mit GPT-4o-mini}: 50x günstiger als GPT-4o,
    für 80\,\% der Aufgaben ausreichend. Nur bei komplexen Fällen
    auf das teure Modell wechseln.
    \item \textbf{Prompt-Caching aktivieren}: Statische Prompt-Teile
    (System-Prompt, Few-Shot-Beispiele) cachen spart 50--90\,\%
    Input-Kosten.
    \item \textbf{Kontext-Länge begrenzen}: Maximal 20\,\% des
    Kontextfensters nutzen. Die Kosten skalieren quadratisch mit der
    Länge, der Nutzen logarithmisch.
    \item \textbf{RAG statt Full-Context}: Statt alle Dokumente in den
    Prompt zu packen, nutze RAG mit Vorauswahl. Das spart Kosten und
    verbessert die Qualität (Lost-in-the-Middle-Effekt).
    \item \textbf{Output-Länge kontrollieren}: Setze \texttt{max\_tokens}
    auf das notwendige Minimum. Output-Tokens sind 3--10x teurer als
    Input-Tokens.
    \item \textbf{Budget-Alarme einrichten}: Setze ein monatliches
    Budget-Limit pro API-Key. Bei Überschreitung schlägt ein Alarm.
    \item \textbf{Modelle nach Aufgaben staffeln}: Einfache Aufgaben
    mit GPT-4o-mini (oder Gemini Flash), komplexe mit GPT-4o (oder
    Claude 4). Ein Semantic Router automatisiert die Entscheidung.
\end{itemize}

% ===========================================================%
% ANTI-PATTERNS
% ============================================================
\section{Anti-Patterns}

\begin{description}[leftmargin=*,style=nextline]
    \item[Maximales Kontextfenster ausreizen] Nur weil ein Modell
    1M Token unterstützt, heisst das nicht, dass du sie nutzen solltest.
    Die Kosten steigen quadratisch, die Qualität leidet (Lost in the
    Middle). Nutze maximal 20\,\% des Fensters.

    \item[Keine Kostenkontrolle] Ohne Budget-Limits und Monitoring
    kann ein Bug in der Prompt-Logik schnell Tausende Euro kosten.
    Ein User, der immer mehr Dokumente an den Prompt hängt, explodiert
    die Kosten.

    \item[Alle Aufgaben mit dem teuersten Modell] GPT-4o für jede
    Anfrage zu nutzen ist die häufigste Kostenfalle. Ein Semantic
    Router spart 60--80\,\% der Kosten ohne Qualitätsverlust.

    \item[Kontext-History ohne Grenzen] Chat-History ohne Begrenzung
    wächst mit jeder Nachricht. Nach 100 Nachrichten ist der Prompt
    zu gross -- und zu teuer. Implementiere immer ein Sliding Window
    oder eine Summary Compression.

    \item[Output-Länge nicht begrenzen] Ohne \texttt{max\_tokens} kann
    ein Modell 10.000 Tokens generieren, wo 200 reichen. Das kostet
    Geld und erhöht die Latenz.

    \item[Caching bei dynamischen Prompts erwarten] Prompt-Caching
    funktioniert nur bei exakt identischen Präfixen. Wenn der System-
    Prompt dynamische Werte enthält, wird nicht gecached.
\end{description}

% ===========================================================%
% PERFORMANCE
% ============================================================
\section{Performance und Optimierung}

\subsection{Multi-Model-Routing}

Der effektivste Kostenhebel: Nicht jedes Modell für jede Aufgabe:

\begin{table}[H]
\centering
\begin{tabularx}{\linewidth}{lXX}
\toprule
\textbf{Aufgabentyp} & \textbf{Modell} & \textbf{Kosten/Query} \\
\midrule
Einfache Klassifikation & GPT-4o-mini & 0,001 \$ \\
Extraktion (strukturiert) & GPT-4o-mini & 0,002 \$ \\
Zusammenfassung (kurz) & GPT-4o-mini & 0,003 \$ \\
Komplexes Reasoning & GPT-4o / Claude 4 & 0,05--0,10 \$ \\
Codegenerierung & DeepSeek-Coder & 0,005 \$ \\
\bottomrule
\end{tabularx}
\caption{Kosten pro Query nach Aufgabentyp}
\end{table}

\subsection{Semantic Router}

\begin{lstlisting}[language=Python, caption={Semantic Router fur Modellauswahl}, label=lst:semantic-router]
def route_query(query: str) -> str:
    """Entscheide, welches Modell fur diese Query optimal ist."""
    prompt = (
        "Klassifiziere die folgende User-Anfrage in eine Kategorie: "
        "'simple' (1-2 Satze, Faktenabfrage, Extraktion), "
        "'complex' (Reasoning, Analyse, Multi-Step), "
        "'creative' (Brainstorming, Textgenerierung).\n\n"
        f"Anfrage: {query}\n\n"
        "Antworte nur mit der Kategorie."
    )

    response = client.chat.completions.create(
        model="gpt-4o-mini",
        messages=[{"role": "user", "content": prompt}],
        max_tokens=10,
        temperature=0.0,
    )
    category = response.choices[0].message.content.strip()

    routing = {
        "simple": "gpt-4o-mini",
        "complex": "gpt-4o",
        "creative": "claude-sonnet-4",
    }
    return routing.get(category, "gpt-4o-mini")

# Nutzung in der API
model = route_query(user_query)
response = call_model(model, user_query)
\end{lstlisting}

\subsection{Kosten-Monitoring}

\begin{lstlisting}[language=Python, caption={Kosten-Tracker fur LLM-API-Calls}, label=lst:cost-tracker]
import time
from collections import defaultdict

class CostTracker:
    """Trackt API-Kosten pro Modell und Tag."""

    def __init__(self, budget_monthly: float = 100.0):
        self.usage: dict[str, list[dict]] = defaultdict(list)
        self.budget = budget_monthly

    def track(self, model: str, input_tokens: int,
              output_tokens: int) -> float:
        pricing = MODEL_PRICING[model]
        cost = (input_tokens / 1_000_000
                * pricing.input_price_per_million
                + output_tokens / 1_000_000
                * pricing.output_price_per_million)
        self.usage[model].append({
            "cost": cost,
            "timestamp": time.time(),
        })
        return cost

    def monthly_cost(self) -> float:
        """Gesamtkosten des aktuellen Monats."""
        month_ago = time.time() - 30 * 24 * 3600
        total = 0.0
        for model, calls in self.usage.items():
            for call in calls:
                if call["timestamp"] > month_ago:
                    total += call["cost"]
        return total

    def check_budget(self) -> str | None:
        """Warnung bei Budget-uberschreitung."""
        used = self.monthly_cost()
        if used > self.budget * 0.8:
            return (
                f"WARNUNG: {used:.2f} $ von {self.budget} $ "
                f"verbraucht ({used/self.budget:.0%})"
            )
        return None
\end{lstlisting}

\begin{praxishinweis}
Kosten-Monitoring ohne Alarme ist sinnlos. Ein Cost-Exploit (Angreifer füllt den Kontext mit 100K Tokens) oder ein fehlerhafter Batch-Job können innerhalb von Minuten das Monatsbudget aufbrauchen. Setze harte Limits pro Request, pro User und pro Tag.
\end{praxishinweis}

% ===========================================================%
% SICHERHEIT
% ============================================================
\section{Sicherheit -- Kontext als Angriffsfläche}

\begin{itemize}[leftmargin=*]
    \item \textbf{Kontext-Injection über lange Prompts}: Ein Angreifer
    kann durch sehr lange Inputs das Context Window füllen und so
    die Kosten in die Höhe treiben (Cost-Exploit). Setze harte
    Limits pro Request.
    \item \textbf{History-Poisoning}: Wenn User die Chat-History
    manipulieren können (z.B. durch vorherige Prompt Injection),
    kann das den Kontext vergiften. Validiere die History vor
    jedem neuen Call.
    \item \textbf{PII im Kontext}: Personenbezogene Daten, die im
    RAG-Kontext oder in der Chat-History enthalten sind, werden
    an die API gesendet. Maskiere PII vor dem Call.
    \item \textbf{Cache-Leakage}: Gecachte Prompts könnten sensible
    Informationen enthalten (Prompt-Engineering-Geheimnisse,
    System-Prompt mit Geschäftslogik). Teile Caches nie zwischen
    verschiedenen Mandanten.
\end{itemize}

% ===========================================================%
% ZUSAMMENFASSUNG
% ===========================================================%
\begin{zusammenfassung}
\begin{itemize}[leftmargin=*]
    \item Kontext kostet Geld -- jeder Token im Prompt wird berechnet.
    Input-Kosten dominieren bei RAG-Systemen (70--90\,\%).
    \item Output-Tokens sind 3--10x teurer als Input-Tokens.
    Setze \texttt{max\_tokens} immer auf das Minimum.
    \item Prompt-Caching spart 50--90\,\% der Input-Kosten bei
    wiederholten System-Prompts.
    \item Lost in the Middle: Wichtige Informationen gehören an den
    Anfang oder das Ende des Kontexts, nicht in die Mitte.
    \item Multi-Model-Routing spart 60--80\,\% der Kosten ohne
    Qualitätsverlust.
    \item Grösseres Kontextfenster ist kein Ersatz für gutes RAG.
    Vorauswahl ist besser als Volltext.
\end{itemize}
\end{zusammenfassung}

% ===========================================================%
% MERKE-BOX
% ===========================================================%
\begin{merke}
\begin{itemize}
    \item \textbf{Kontext ist teuer}: Jeder Token kostet Geld.
    \item \textbf{Output ist teurer als Input}: 3--10x Faktor.
    \item \textbf{Caching ist der effektivste Sparhebel}.
    \item \textbf{Grosser Kontext != gute Antwort}: Lost in the Middle.
    \item \textbf{Modell-Routing spart 80\,\% Kosten}.
\end{itemize}
\end{merke}

% ===========================================================%
% PRAXISPROJEKT
% ===========================================================%
\section{Praxisprojekt -- Kostenoptimierungs-Dashboard}

\textbf{Ziel:} Baue ein Dashboard, das die LLM-Kosten deiner Anwendung
in Echtzeit verfolgt und Optimierungsvorschläge macht.

\textbf{Aufgaben:}
\begin{enumerate}[leftmargin=*]
    \item Implementiere einen CostTracker, der jeden API-Call mit
    Modell, Token-Zahlen und Kosten protokolliert
    \item Baue ein Semantic Router-System, das Queries automatisch
    an das günstigste ausreichende Modell routet
    \item Erstelle einen Report, der zeigt: Kosten pro Tag, pro Modell,
    pro User, Kosten-Entwicklung über Zeit
    \item Füge Budget-Alarme hinzu (80\,\%-, 90\,\%-, 100\,\%-Schwellen)
    \item Optimiere ein existierendes System um mindestens 50\,\%
    Kostenreduktion durch: kleineres Modell + Caching + Output-Limits
\end{enumerate}

\textbf{Testdaten:} Simuliere 1.000 API-Calls mit verschiedenen Modellen
und Prompt-Längen. Zeige die Kostenunterschiede der Optimierungsstufen.

\textbf{Erfolgskriterium:} Das Dashboard zeigt korrekte Tageskosten,
warnt bei Budget-Überschreitung und macht automatisch
Optimierungsvorschläge.

% ===========================================================%
% WEITERFÜHRENDE RESSOURCEN
% ============================================================
\section{Weiterführende Ressourcen}

\begin{itemize}[leftmargin=*]
    \item \textbf{OpenAI Tokenizer}: \href{https://platform.openai.com/tokenizer}{platform.openai.com/tokenizer}
    -- Online-Tool zum Zählen und Analysieren von Tokens
    \item \textbf{tiktoken (GitHub)}: \href{https://github.com/openai/tiktoken}{github.com/openai/tiktoken}
    -- Pythons Tokenizer für OpenAI-Modelle
    \item \textbf{Anthropic Prompt Caching Guide}: \href{https://docs.anthropic.com/en/docs/build-with-claude/prompt-caching}{docs.anthropic.com}
    \item \textbf{OpenAI Pricing}: \href{https://openai.com/pricing}{openai.com/pricing}
    \item \textbf{Lost in the Middle Paper}: Liu et al. (2024):
    \glqq Lost in the Middle: How Language Models Use Long Contexts\grqq
    \item \textbf{Semantic Router (GitHub)}: \href{https://github.com/aurelio-labs/semantic-router}{github.com/aurelio-labs/semantic-router}
    \item \textbf{Langfuse Cost Tracking}: \href{https://langfuse.com/docs/cost-tracking}{langfuse.com/docs/cost-tracking}
\end{itemize}

\textbf{Nächster Schritt:} Jetzt kennst du die Kosten- und Kontext-Fakten.
Im nächsten Kapitel lernst du die OpenAI Platform im Detail kennen -- von
der API-Nutzung bis zu Best Practices für Production.
